% ROZDZIAŁ 23 — MIARA HAARA
\chapter{Miara Haara na grupach lokalnie zwartych}
\label{ch:haar-23}

\section{Grupy topologiczne, translacje i miary niezmiennicze}
\begin{definicja}[Grupa topologiczna]
\label{def:topological-group-23}
Grupę $G$, wyposażoną w topologię Hausdorffa, nazywamy
\emph{grupą topologiczną}, gdy odwzorowania $(x,y)\mapsto xy$
oraz $x\mapsto x^{-1}$ są ciągłe. Element neutralny oznaczamy $e$.
Grupa jest \emph{lokalnie zwarta}, gdy każdy punkt ma otoczenie
o zwartej domkniętej powłoce. Piszemy $xA=\{xa:a\in A\}$
i $Ax=\{ax:a\in A\}$.
\end{definicja}

\begin{uwaga}[Dlaczego lokalna zwartość?]
W przestrzeni $\mathbb R^d$ miara Lebesgue'a jest niezmiennicza
względem przesunięć. Dla grup nieprzemiennych należy rozróżniać
translacje \emph{lewe} i \emph{prawe}. Lokalna zwartość gwarantuje,
że zbiory zwarte można pokrywać skończenie wieloma przesunięciami
małego otwartego otoczenia $e$; na tej własności oprzemy konstrukcję.
\end{uwaga}

\begin{definicja}[Lewostronna i prawostronna miara Haara]
\label{def:haar-23}
Niech $G$ będzie lokalnie zwartą grupą Hausdorffa.
\emph{Lewostronną miarą Haara} nazywamy niezerową miarę Radona
$\mu$ na $\mathcal B(G)$, w sensie definicji~\ref{def:radon-measure-20},
dla której $\mu(gE)=\mu(E)$ dla wszystkich $g\in G$
i $E\in\mathcal B(G)$. Analogicznie definiujemy
\emph{prawostronną miarę Haara} przez warunek
$\mu(Eg)=\mu(E)$. Zakres wymaganej regularności
przypominamy poniżej ze względu na brak założeń przeliczalności.
\end{definicja}

\begin{twierdzenie}[Haara: istnienie i jednoznaczność]
\label{thm:haar-main-23}
Na każdej lokalnie zwartej grupie Hausdorffa istnieje
lewostronna miara Haara. Każde dwie takie miary różnią się
dodatnim czynnikiem stałym. To samo dotyczy prawostronnych miar
Haara.
\end{twierdzenie}

\begin{uwaga}[Zakres twierdzenia i regularność]
W twierdzeniu~\ref{thm:haar-main-23} \emph{nie} zakładamy ani
metryzowalności, ani $\sigma$-zwartości, ani drugiej przeliczalności.
W ogólnej przestrzeni lokalnie zwartej Hausdorffa rozróżniamy
trzy warunki dla miary borelowskiej $\mu$:
\begin{enumerate}[label=(\roman*)]
\item $\mu(K)<\infty$ dla każdego kompaktu $K$;
\item $\mu(E)=\inf\{\mu(U):E\subseteq U,\ U\text{ otwarty}\}$
      dla każdego zbioru borelowskiego $E$;
\item $\mu(U)=\sup\{\mu(K):K\subseteq U,\ K\text{ zwarty}\}$
      dla każdego zbioru otwartego $U$.
\end{enumerate}
To właśnie te trzy własności przyjmujemy w definicji miary Radona
na dowolnej grupie lokalnie zwartej. Nie wymagamy regularności
wewnętrznej na \emph{każdym} zbiorze borelowskim: bez
$\sigma$-zwartości grupy jest to założenie zbyt silne.
Dla borelowskich zbiorów o mierze $\sigma$-skończonej
regularność wewnętrzna wyniknie z (i)--(iii).
Wyjaśnimy różnicę na konkretnym przykładzie na końcu rozdziału.
\end{uwaga}

\section{Lematy topologiczne poprzedzające konstrukcję}
\begin{lemat}[Małe otoczenia, zwarte rozdrobnienia i pokrycia]
\label{lem:haar-neigh-23}
Niech $G$ będzie lokalnie zwartą grupą Hausdorffa.
\begin{enumerate}[label=(\roman*)]
\item Dla każdego otoczenia $O$ elementu $e$ istnieje symetryczne,
  relatywnie zwarte otoczenie otwarte $V$ takie, że
  $V^{-1}V\subseteq O$ oraz $\overline V\subseteq O$.
  Relatywnie zwarte otoczenia otwarte
  tworzą bazę otoczeń $e$.
\item Dla każdego kompaktu $K\subseteq G$ i otwartego otoczenia
  $V$ punktu $e$ istnieje skończone pokrycie $K$
  przez lewe przesunięcia $V$.
\item Jeśli $K,L$ są rozłącznymi kompaktami, to istnieje otoczenie
  $V$ punktu $e$ takie, że żaden zbiór $xV$, $x\in G$,
  nie przecina jednocześnie $K$ i $L$.
\item Jeśli $K\subseteq\bigcup_{j=1}^r O_j$ jest zwarty,
  gdzie $O_j$ są otwarte, to istnieją kompakty
  $K_j\subseteq O_j$ spełniające $K=\bigcup_{j=1}^rK_j$.
\end{enumerate}
\end{lemat}
\begin{proof}
Dla (i) z ciągłości mnożenia wybieramy symetryczne
otoczenie otwarte $W$ takie, że $W^2\subseteq O$.
Z lokalnej zwartości istnieje zwarty $C$ z
$e\in\operatorname{int}C$. Kładziemy
$V=W\cap\operatorname{int}C\cap(\operatorname{int}C)^{-1}$.
Wówczas $V$ jest otwarty i symetryczny,
$\overline V\subseteq C$ oraz
$V^{-1}V\subseteq W^2\subseteq O$.
Ponadto $\overline W\subseteq W^2$: jeśli $x\in\overline W$,
to otwarty zbiór $xW$ przecina $W$, a więc
$x=uw^{-1}$ dla pewnych $u,w\in W$.
Stąd także $\overline V\subseteq O$.
Dobór $O$ wewnątrz dowolnego zadanego otoczenia
daje bazę relatywnie zwartych otoczeń.

W (ii) rodzina $(xV)_{x\in K}$ jest pokryciem otwartym
kompaktu $K$, skąd istnieje podpokrycie skończone.
Dla (iii) kompakt $K^{-1}L$ nie zawiera $e$;
stosujemy (i) do otwartego otoczenia
$G\setminus K^{-1}L$. Gdyby $xv\in K$ i $xw\in L$
przy $v,w\in V$, otrzymalibyśmy
$v^{-1}w\in K^{-1}L\cap V^{-1}V$, sprzeczność.

W (iv) wykorzystamy następujący fakt o każdej przestrzeni
lokalnie zwartej Hausdorffa $X$: jeśli $x\in O$ i $O$
jest otwarty, to istnieje otwarty $W\ni x$ ze
zwartym $\overline W\subseteq O$. Dla dowodu wybieramy
zwartą okolicę $C$ punktu $x$, z $x\in\operatorname{int}C$.
Zbiór $F=C\setminus O$ jest zwartym podzbiorem
przestrzeni Hausdorffa $C$. Dla każdego $y\in F$
rozdzielamy $x$ i $y$ rozłącznymi otoczeniami w $C$.
Zwartość $F$ pozwala wybrać skończenie wiele z nich;
przecięcie otoczeń $x$ daje otwarty $Z\subseteq C$
taki, że $x\in Z$ i $\overline Z^{\,C}\cap F=\varnothing$.
Bierzemy otwarty $\widetilde Z\subseteq X$ z
$\widetilde Z\cap C=Z$ i definiujemy
$W=\widetilde Z\cap\operatorname{int}C$.
Wtedy $\overline W\subseteq C\cap\overline Z^{\,C}
\subseteq O$ oraz $\overline W$ jest zwarty.

Dla każdego $x\in K$ wybieramy zatem indeks $j(x)$
oraz takie otoczenie $W_x$, aby
$x\in W_x\subseteq\overline W_x\subseteq O_{j(x)}$.
Wybierzmy $x_1,\dots,x_s$ tak, aby zbiory $W_{x_i}$
pokrywały $K$, i połóżmy
$K_j=K\cap\bigcup_{i:\,j(x_i)=j}\overline W_{x_i}$.
Wówczas każde $K_j$ jest zwarte, zawiera się w $O_j$,
a ich suma zbiorowa jest równa $K$.
\end{proof}

\section{Zwarta przestrzeń danych bez założenia przeliczalności}
\begin{lemat}[Zwartość dowolnego produktu przedziałów]
\label{lem:haar-product-compact-23}
Niech $I$ będzie dowolnym zbiorem indeksów, a
$a_i\in[0,\infty)$ dla $i\in I$. W topologii produktowej
przestrzeń $P=\prod_{i\in I}[0,a_i]$ jest zwarta.
\end{lemat}
\begin{proof}
Dla zupełności podajemy argument ultrafiltrowy, nie wymagający
uprzedniego użycia twierdzenia Tichonowa. Korzystamy jedynie
z lematu Kuratowskiego--Zorna, czyli ze standardowej postaci
aksjomatu wyboru. Jeśli rodzina zbiorów ma własność skończonych
przecięć, rodzina wszystkich filtrów ją zawierających jest
niepusta. Każdy łańcuch tych filtrów ma ograniczenie górne,
którym jest ich suma zbiorowa. Lemat Zorna dostarcza filtru
maksymalnego $\mathcal U$. Jest to ultrafiltr: gdy $A\notin\mathcal U$, maksymalność oznacza, że
rodzina $\mathcal U\cup\{A\}$ nie może być bazą filtru
właściwego. Istnieje więc $B\in\mathcal U$ z
$B\cap A=\varnothing$, skąd $X\setminus A\in\mathcal U$
($X$ jest tutaj zbiorem, na którym zdefiniowano ultrafiltr).

Aby dowieść zwartości $P$, niech $(F_\alpha)_{\alpha\in J}$
będzie rodziną domkniętych podzbiorów $P$ o własności
skończonych przecięć. Rozszerzamy ją do ultrafiltru
$\mathcal U$ na $P$. Dla każdego $i\in I$ obraz
$\mathcal U$ przez rzutowanie $\pi_i:P\to[0,a_i]$
jest ultrafiltrem na zwartym przedziale. Zwartość przedziału
i własność skończonych przecięć zapewniają, że
$C_i:=\bigcap_{A\in\pi_i(\mathcal U)}\overline A$ jest niepusty.
Wybieramy $t_i=\min C_i$ (nie trzeba dokonywać dodatkowego
wyboru punktów). Każde otwarte otoczenie $W$ punktu $t_i$
należy do $\pi_i(\mathcal U)$: w przeciwnym razie
$[0,a_i]\setminus W$ należałoby do ultrafiltru,
co wykluczałoby $t_i\in C_i$.

Punkt $t=(t_i)_{i\in I}$ należy do domknięcia każdego
zbioru z $\mathcal U$. Istotnie, podstawowe otoczenie
$t$ ogranicza tylko skończenie wiele współrzędnych,
więc należy do $\mathcal U$ i musi przecinać każdy
jego element. W szczególności $t\in F_\alpha$ dla
każdego $\alpha$, ponieważ $F_\alpha$ jest domknięty
i należy do $\mathcal U$. Wobec tego każda rodzina
zbiorów domkniętych o własności skończonych przecięć
ma niepuste przecięcie, co dowodzi zwartości $P$.
\end{proof}

\section{Konstrukcja miary przez liczby pokrycia}
Niech odtąd $G$ będzie \emph{dowolną} lokalnie zwartą
grupą Hausdorffa. Wybieramy kompakt $K_0$ będący otoczeniem $e$
oraz otwarte otoczenie $V_0$ elementu $e$,
spełniające $V_0\subseteq K_0$. Przez $\mathcal N$
oznaczamy rodzinę wszystkich relatywnie zwartych otwartych
otoczeń $e$; jest ona skierowana przez odwrotną inkluzję.
Dla zwartego $K$ i $V\in\mathcal N$ definiujemy
\begin{equation}\label{eq:haar-cover-number-23}
 N(K,V)=\min\left\{r\in\mathbb N_0:\
 K\subseteq\bigcup_{j=1}^{r}x_jV
 \text{ dla pewnych }x_1,\dots,x_r\in G\right\}.
\end{equation}
W szczególności $N(\varnothing,V)=0$, a dla
$K\ne\varnothing$ zachodzi $1\le N(K,V)<\infty$.

\begin{lemat}[Jednolite oszacowania ilorazów]
\label{lem:haar-ratio-23}
Niech $C_K=N(K,V_0)$. Dla każdego $V\in\mathcal N$ mamy
\[
 0\le q_V(K):=\frac{N(K,V)}{N(K_0,V)}\le C_K,
 \qquad q_V(K_0)=1.
\]
Ponadto $q_V(gK)=q_V(K)$, funkcje $q_V$ są monotoniczne
i podaddytywne dla sum zwartych, zaś dla rozłącznych
kompaktów $K,L$ istnieje $W\in\mathcal N$ takie, że
$q_V(K\cup L)=q_V(K)+q_V(L)$ dla każdego $V\subseteq W$.
\end{lemat}
\begin{proof}
Zbiory $K$ można pokryć przez $C_K$ przesunięć $V_0$,
a każde $xV_0\subseteq xK_0$. Pokrywając $K_0$
przez $N(K_0,V)$ przesunięć $V$ i następnie przesuwając
to pokrycie, otrzymujemy $N(K,V)\le C_KN(K_0,V)$.
Niezmienniczość i podaddytywność wynikają z przesuwania
oraz łączenia skończonych pokryć.

Dla rozłącznych $K,L$ wybierzmy $W$ według
lematu~\ref{lem:haar-neigh-23}(iii). Gdy $V\subseteq W$,
żadne przesunięcie $V$ nie przecina zarazem $K$ i $L$.
Dowolne minimalne pokrycie $K\cup L$ można zatem
podzielić na dwie rozłączne rodziny, z których pierwsza
pokrywa $K$, a druga $L$. Stąd
$N(K\cup L,V)\ge N(K,V)+N(L,V)$.
Nierówność odwrotna wynika z łączenia pokryć.
\end{proof}

\begin{lemat}[Zawartość niezmiennicza na kompaktach]
\label{lem:haar-content-23}
Istnieje skończona funkcja nieujemna $q$ na rodzinie
wszystkich kompaktów $K\subseteq G$, taka że
\begin{enumerate}[label=(\roman*)]
\item $q(\varnothing)=0$, $q(K_0)=1$ oraz $q(gK)=q(K)$;
\item $K\subseteq L$ implikuje $q(K)\le q(L)$;
\item $q(K\cup L)\le q(K)+q(L)$;
\item jeśli $K\cap L=\varnothing$, to
      $q(K\cup L)=q(K)+q(L)$.
\end{enumerate}
\end{lemat}
\begin{proof}
Niech $\mathcal K$ będzie zbiorem wszystkich zwartych
podzbiorów $G$. Na mocy
lematu~\ref{lem:haar-product-compact-23} produkt
$P=\prod_{K\in\mathcal K}[0,C_K]$ jest zwarty,
bez względu na moc $\mathcal K$.
Każde $q_V=(q_V(K))_{K\in\mathcal K}$ należy do $P$.
Dla $W\in\mathcal N$ połóżmy
$F_W=\overline{\{q_V:V\in\mathcal N,\ V\subseteq W\}}$.
Jest to niepusty zbiór domknięty. Dla skończonej
rodziny $W_1,\dots,W_r$ mamy
$F_{W_1\cap\dots\cap W_r}\subseteq\bigcap_{j=1}^rF_{W_j}$.
Zwartość $P$ pozwala wybrać
$q\in\bigcap_{W\in\mathcal N}F_W$.

Własności (i)--(iii) opisują domknięte warunki na
skończenie wielu współrzędnych, spełnione przez
każde $q_V$, więc także przez $q$.
Jeżeli $K\cap L=\varnothing$, z poprzedniego lematu
mamy otoczenie $W$, poniżej którego wszystkie $q_V$
spełniają równość addytywności dla $K,L$.
Jej zbiór rozwiązań jest domknięty i zawiera $F_W$,
a więc również $q$. To dowodzi (iv).
\end{proof}

\begin{lemat}[Od zawartości na kompaktach do miary borelowskiej]
\label{lem:haar-content-to-measure-23}
Niech $X$ będzie lokalnie zwartą przestrzenią Hausdorffa,
a $q$ skończoną funkcją nieujemną na kompaktach,
spełniającą (i)--(iv) z poprzedniego lematu z pominięciem
niezmienniczości i normalizacji. Dla otwartego $U$
i dowolnego $A\subseteq X$ definiujemy
\begin{equation}\label{eq:haar-from-content-23}
 t(U)=\sup_{K\subseteq U,\ K\ \mathrm{zwarty}}q(K),
 \qquad
 m^*(A)=\inf_{U\supseteq A,\ U\ \mathrm{otwarty}}t(U).
\end{equation}
Wówczas $m^*$ jest miarą zewnętrzną,
wszystkie zbiory borelowskie są mierzalne w sensie
Carathéodory'ego, a $\mu=m^*|_{\mathcal B(X)}$
spełnia warunki (i)--(iii) regularności podane
na początku rozdziału. Ponadto $\mu(K)\ge q(K)$.
Jeśli homeomorfizm $T:X\to X$ zachowuje $q$ na
kompaktach, to $\mu(TA)=\mu(A)$ dla każdego
zbioru borelowskiego $A$.
\end{lemat}
\begin{proof}
Wprost z definicji $m^*(U)=t(U)$, gdy $U$ jest otwarty.
Najpierw dowodzimy podaddytywności $t$ na pokryciach
otwartych. Jeśli $K\subseteq\bigcup_{j\ge1}U_j$ jest
zwarty, możemy wybrać skończone podpokrycie.
Z lematu~\ref{lem:haar-neigh-23}(iv) (którego dowód
jest topologiczny i działa w każdej przestrzeni lokalnie
zwartej Hausdorffa) dostajemy kompakty
$K_j\subseteq U_j$ pokrywające $K$. Podaddytywność $q$
daje $q(K)\le\sum_jt(U_j)$. Biorąc supremum po $K$,
uzyskujemy
$t(\bigcup_{j\ge1}U_j)\le\sum_jt(U_j)$.
Aby uzyskać przeliczalną podaddytywność $m^*$,
dla zbiorów $A_j$ o skończonej sumie $\sum_jm^*(A_j)$
i $\varepsilon>0$ wybieramy otwarte $U_j\supseteq A_j$
z $t(U_j)<m^*(A_j)+\varepsilon2^{-j}$.
Wtedy $m^*(\bigcup_jA_j)\le\sum_jm^*(A_j)+\varepsilon$;
pozostałe własności miary zewnętrznej są natychmiastowe.

Najważniejszym krokiem jest mierzalność \emph{otwartych}
zbiorów bez odwoływania się do jakiejkolwiek metryki.
Ustalmy otwarty $O\subseteq X$, dowolny zbiór $A$
o $m^*(A)<\infty$ i $\varepsilon>0$.
Wybieramy otwarty $V\supseteq A$ z
$t(V)<m^*(A)+\varepsilon$. Ponieważ
$t(V\cap O)\le t(V)<\infty$, istnieje zwarty
$K\subseteq V\cap O$ taki, że
$q(K)>t(V\cap O)-\varepsilon$; jeśli supremum
jest zerowe, bierzemy $K=\varnothing$.
Dla dowolnego kompaktu $L\subseteq V\setminus K$
addytywność na rozłącznych kompaktach daje
$q(K)+q(L)=q(K\cup L)\le t(V)$.
Stąd $t(V\setminus K)\le t(V)-q(K)$.
Ponieważ $A\setminus O\subseteq V\setminus K$,
a $A\cap O\subseteq V\cap O$, dostajemy
\[
 m^*(A\setminus O)+m^*(A\cap O)
 \le t(V\setminus K)+t(V\cap O)
 \le t(V)+\varepsilon<m^*(A)+2\varepsilon.
\]
Przy $m^*(A)=\infty$ potrzebna nierówność
jest automatyczna. Po przejściu z $\varepsilon$ do
zera warunek Carathéodory'ego jest spełniony.
Z twierdzenia~\ref{thm:caratheodory} wynika,
że $\mu$ jest miarą na $\mathcal B(X)$.

Regularność zewnętrzna jest zawarta w definicji $m^*$.
Dla kompaktu $K$ wybieramy otwarty $W\supseteq K$
o zwartej domkniętej powłoce $C$. Wówczas
$\mu(K)\le t(W)\le q(C)<\infty$.
Dla każdego otwartego $U\supseteq K$ mamy
$t(U)\ge q(K)$, więc $\mu(K)\ge q(K)$.
Dla otwartego $O$ zachodzi
$\mu(O)=t(O)=\sup_{K\subseteq O}q(K)
\le\sup_{K\subseteq O}\mu(K)\le\mu(O)$,
co dowodzi regularności wewnętrznej na otwartych.

Na koniec, jeśli $T$ jest wskazanym homeomorfizmem,
to $T$ przekształca bijektywnie kompakty zawarte w $U$
na kompakty zawarte w $TU$, skąd
$t(TU)=t(U)$ i $m^*(TA)=m^*(A)$.
\end{proof}

\begin{proof}[Dowód istnienia w twierdzeniu~\ref{thm:haar-main-23}
bez założeń przeliczalności]
Z lematu~\ref{lem:haar-content-23} otrzymujemy
zawartość $q$. Lemat~\ref{lem:haar-content-to-measure-23}
daje borelowską miarę Radona $\mu$ niezmienniczą
względem wszystkich lewych translacji.
Ponieważ $\mu(K_0)\ge q(K_0)=1$, jest niezerowa.
Jej skończoność na kompaktach wynika z konstrukcji.
Dla miary prawej stosujemy inwersję:
$\mu_R(E)=\mu(E^{-1})$.
Jest to miara Radona, bo inwersja jest homeomorfizmem,
a $\mu_R(Eg)=\mu(g^{-1}E^{-1})=\mu(E^{-1})$.
Nie było potrzebne żadne założenie o przeliczalnej bazie
otoczeń $e$ ani o $\sigma$-zwartości $G$.
\end{proof}

\begin{wniosek}[Dodatniość miary zbiorów otwartych]
\label{cor:haar-open-positive-23}
Jeśli $\mu$ jest niezerową lewą miarą Haara
na lokalnie zwartej grupie Hausdorffa $G$, to
$\mu(U)>0$ dla każdego niepustego otwartego $U\subseteq G$.
\end{wniosek}
\begin{proof}
Jeśli $\mu(U)=0$ dla pewnego niepustego $U$, wybierzmy
$a\in U$. Dla każdego kompaktu $K\subseteq G$ skończenie
wiele zbiorów $xa^{-1}U$ pokrywa $K$, ponieważ każdy
z nich jest przesunięciem otwartego otoczenia $e$.
Z niezmienniczości wynika $\mu(K)=0$.
Regularność wewnętrzna na otwartych daje teraz
$\mu(G)=0$, wbrew niezerowości miary.
\end{proof}

\section{Jednoznaczność: argument z funkcją przybliżającą jedynkę}
\begin{lemat}[Jednostajna ciągłość translacji funkcji]
\label{lem:haar-uniform-23}
Jeśli $f\in C_c(G)$, to dla każdego $\varepsilon>0$
istnieje otoczenie $V$ elementu $e$ takie, że
$|f(yz)-f(y)|<\varepsilon$ dla wszystkich
$y\in G$ i $z\in V$.
\end{lemat}
\begin{proof}
Niech $K=\operatorname{supp}f$ i wybierzmy symetryczne
relatywnie zwarte otoczenie $C$ elementu $e$.
Zbiór $M=K\overline C$ jest zwarty. Ciągłość
$(y,z)\mapsto f(yz)-f(y)$ w każdym $(y,e)$,
$y\in M$, pozwala wybrać otoczenia $W_y$ punktu $y$
i $V_y$ elementu $e$ takie, że
$|f(uz)-f(u)|<\varepsilon$ dla $u\in W_y$, $z\in V_y$.
Ze zwartego $M$ wybieramy skończone pokrycie
$W_{y_1},\dots,W_{y_r}$; połóżmy
$V=C\cap V_{y_1}\cap\cdots\cap V_{y_r}$,
zmniejszając je w razie potrzeby do symetrycznego
otoczenia. Dla $u\in M$, $z\in V$ teza wynika
z konstrukcji. Jeśli $u\notin M$, to $f(u)=0$,
a $uz\in K$ prowadziłoby do
$u\in K V^{-1}\subseteq K\overline C=M$.
Zatem również $f(uz)=0$.
\end{proof}

\begin{lemat}[Fubini dla ciągłych jąder o zwartym nośniku]
\label{lem:haar-fubini-compact-23}
Niech $X,Y$ będą lokalnie zwartymi przestrzeniami Hausdorffa,
a $\alpha,\beta$ miarami borelowskimi skończonymi na
kompaktach. Jeśli $H\in C_c(X\times Y)$, to obie całki iterowane
\[
 \int_X\!\left(\int_Y H(x,y)\,\mathrm d\beta(y)\right)
 \mathrm d\alpha(x),\qquad
 \int_Y\!\left(\int_X H(x,y)\,\mathrm d\alpha(x)\right)
 \mathrm d\beta(y)
\]
są dobrze określone, skończone i równe, nawet gdy
$X,Y$ nie są metryzowalne ani $\alpha,\beta$
nie są $\sigma$-skończone.
\end{lemat}
\begin{proof}
Rzuty zwartego nośnika $H$ są kompaktami $K\subseteq X$
i $L\subseteq Y$; całkowanie ograniczamy do $K\times L$.
Na iloczynie tych zwartych przestrzeni każdą funkcję
ciągłą przybliża się jednostajnie przez sumy
$\sum_{j=1}^Nc_j a_j(x)b_j(y)$, gdzie
$a_j\in C(K)$, $b_j\in C(L)$. Oto bezpośrednie
uzasadnienie tej aproksymacji. Dla danego $\varepsilon>0$
z ciągłości $H$ i zwartości iloczynu można wybrać
skończone otwarte pokrycia $(U_i)$ zbioru $K$ oraz
$(V_j)$ zbioru $L$ takie, że oscylacja $H$ na każdym
$U_i\times V_j$ jest mniejsza niż $\varepsilon$.
Istnienie tych pokryć sprawdzamy bez metryzowalności:
dla ustalonego $x\in K$ i każdego $y\in L$ ciągłość
$H$ daje otoczenia $U_{x,y}\ni x$, $V_{x,y}\ni y$
takie, że $|H(x',y')-H(x,y)|<\varepsilon/4$
dla $(x',y')\in U_{x,y}\times V_{x,y}$.
Zwartość $L$ pozwala wybrać skończenie wiele
$V_{x,y}$ pokrywających $L$; przecięcie
odpowiednich $U_{x,y}$ daje otoczenie $U_x$,
na którym dla dowolnych $x',x''\in U_x$ i $y\in L$
mamy $|H(x',y)-H(x'',y)|<\varepsilon/2$.
Analogicznie dobieramy otoczenia $V_y$ z oszacowaniem
$|H(x,y')-H(x,y'')|<\varepsilon/2$ jednostajnym
względem $x\in K$. Skończone podpokrycia
$(U_i)$, $(V_j)$ dają żądane oszacowanie oscylacji
na prostokątach $U_i\times V_j$.
Na zwartych przestrzeniach
Hausdorffa istnieją ciągłe podziały jedności
$(a_i)$, $(b_j)$ podporządkowane tym pokryciom.
Wybierając $c_{ij}$ jako wartość $H$ w dowolnym
punkcie niepustego $U_i\times V_j$, otrzymujemy
\[
 \left\|H-\sum_{i,j}c_{ij}a_i\otimes b_j\right\|_\infty
 \le\varepsilon.
\]
Dla sum elementarnych obie całki iterowane są równe
z liniowości. Skończoność $\alpha(K)$ i $\beta(L)$
pozwala przejść do granicy jednostajnej, także dla
całek funkcji zespolonych. Dowodzi to tezy.
\end{proof}

\begin{twierdzenie}[Jednoznaczność miary Haara]
\label{thm:haar-unique-23}
Jeśli $\mu,\nu$ są dwiema lewostronnymi miarami Haara
na lokalnie zwartej grupie Hausdorffa $G$, to istnieje
$c>0$ takie, że $\nu=c\mu$ na $\mathcal B(G)$.
\end{twierdzenie}
\begin{proof}
Niech $f\in C_c(G)$, a $\varphi\in C_c(G)$ będzie
nieujemna, niezerowa i spełniała
$\int_G\varphi(z)\,\mathrm d\mu(z)=1$.
Funkcja $(x,y)\mapsto f(x)\varphi(y^{-1}x)$
ma nośnik zawarty w zwartym iloczynie
$\operatorname{supp}f\times
(\operatorname{supp}f)(\operatorname{supp}\varphi)^{-1}$.
Stosujemy zatem lemat~\ref{lem:haar-fubini-compact-23},
który nie wymaga $\sigma$-skończoności miary na $G$.
Otrzymujemy
\begin{align*}
 I&=\int_G\!\int_G f(x)\varphi(y^{-1}x)
       \,\mathrm d\nu(y)\,\mathrm d\mu(x)\\
  &=c_\varphi\int_Gf(x)\,\mathrm d\mu(x),\qquad
 c_\varphi:=\int_G\varphi(z^{-1})\,\mathrm d\nu(z),
\end{align*}
bo po podstawieniu $y=xz$ wewnętrzna całka nie
zależy od $x$ (lewostronna niezmienniczość $\nu$).
Zamieniając kolejność całkowania i podstawiając
$x=yz$ (lewostronna niezmienniczość $\mu$), dostajemy
\begin{equation}\label{eq:haar-uniqueness-23}
 c_\varphi\int_G f(x)\,\mathrm d\mu(x)
 =\int_G\!\int_G f(yz)\varphi(z)
       \,\mathrm d\mu(z)\,\mathrm d\nu(y).
\end{equation}
Wybierzmy funkcje $\varphi$ o nośnikach zawartych
w coraz mniejszych otoczeniach $e$, każdorazowo
unormowane przez $\mu(\varphi)=1$. Są dostępne,
ponieważ każda niepusta otwarta część grupy ma
dodatnią miarę Haara: w przeciwnym razie jej
translacje pokryłyby skończenie zwarty zbiór
o dodatniej mierze.
Na mocy lematu~\ref{lem:haar-uniform-23}
wewnętrzna całka po prawej stronie
\eqref{eq:haar-uniqueness-23} zbiega jednostajnie
na odpowiednim zwartym zbiorze do $f(y)$,
a poza nim obie funkcje znikają.
Zatem prawa strona dąży do $\int_Gf(y)\,\mathrm d\nu(y)$.
Wybierając ustaloną funkcję $f_0\in C_c(G)$,
$f_0\ge0$ i $\mu(f_0)>0$, otrzymujemy
$c_\varphi\to c:=\nu(f_0)/\mu(f_0)>0$.
Wstawienie dowolnego $f$ do
\eqref{eq:haar-uniqueness-23} daje
$\nu(f)=c\mu(f)$. Pozostaje wyjaśnić, dlaczego równość całek względem
$\mu$ i $\nu$ dla wszystkich funkcji z $C_c(G)$ oznacza
równość miar również bez założenia metryzowalności $G$.
Dla każdego zwartego $K\subseteq U$, gdzie $U$ jest otwarty,
istnieje funkcja $h\in C_c(G)$ taka, że
$0\le h\le1$, $h|_K=1$ i $\operatorname{supp}h\subseteq U$
(lokalnie zwarta przestrzeń Hausdorffa ma takie funkcje
odcinające). Wówczas $\mu(K)\le\int_Gh(x)\,\mathrm d\mu(x)
=\int_Gh(x)\,\mathrm d\nu(x)\le\nu(U)$. Z regularności
wewnętrznej $\mu(U)\le\nu(U)$. Po zamianie miar
otrzymujemy równość na otwartych, a regularność zewnętrzna
daje równość na wszystkich zbiorach borelowskich.
Stąd $\nu=c\mu$, co kończy dowód.
\end{proof}

\begin{przyklad}[Dlaczego nie wymagamy regularności wewnętrznej wszędzie]
\label{ex:haar-non-sigma-compact-23}
Niech $D$ będzie nieprzeliczalną grupą dyskretną,
a $G=\mathbb R\times D$ grupą z topologią produktową.
Rozważmy domknięty zbiór $Y=\{0\}\times D$.
Każdy kompakt $K\subseteq Y$ jest skończony, ponieważ
rzutowanie $K$ na dyskretną przestrzeń $D$ daje zbiór
zwarty, a więc skończony. Punkty mają miarę Haara zero:
na otwartej podgrupie $\mathbb R\times\{e\}$ miara Haara
jest dodatnią wielokrotnością miary Lebesgue'a, zaś
translacje przenoszą ten wniosek na pozostałe składowe.
Wobec tego $\mu(K)=0$ dla każdego kompaktu $K\subseteq Y$.

Z drugiej strony, jeśli $U\supseteq Y$ jest otwarty,
to $U\cap(\mathbb R\times\{d\})$ jest niepustym
zbiorem otwartym dla każdego $d\in D$, a więc ma
ściśle dodatnią miarę Haara na tej składowej.
Składowe są parami rozłączne i jest ich nieprzeliczalnie
wiele, zatem $\mu(U)=\infty$: gdyby $\mu(U)$ było
skończone, najwyżej przeliczalnie wiele rozłącznych
zbiorów mierzalnych mogłoby mieć dodatnią miarę.
Regularność zewnętrzna daje $\mu(Y)=\infty$, chociaż
$\sup_{K\subseteq Y,\,K\text{ zwarty}}\mu(K)=0$.
Takiego zjawiska nie ma dla zbiorów o mierze
$\sigma$-skończonej, zgodnie z lematem
\ref{lem:radon-sigmafinite-inner-20}.
\end{przyklad}

\section{Prawe miary Haara, funkcja modularna i unimodularność}
\begin{twierdzenie}[Prawa miara Haara i modularność]
\label{thm:haar-modular-23}
Niech $\mu$ będzie lewostronną miarą Haara.
Wówczas $\mu_R(E):=\mu(E^{-1})$ jest prawostronną
miarą Haara. Istnieje dokładnie jeden ciągły
homomorfizm $c:G\to(0,\infty)$ taki, że
\begin{equation}\label{eq:haar-modular-23}
 \mu(Eg)=c(g)\mu(E)\qquad(E\in\mathcal B(G),\ g\in G).
\end{equation}
Grupę nazywamy \emph{unimodularną}, gdy $c\equiv1$.
\end{twierdzenie}
\begin{proof}
Inwersja jest homeomorfizmem, toteż $\mu_R$ jest
miarą Radona. Ponadto
$\mu_R(Eg)=\mu((Eg)^{-1})=\mu(g^{-1}E^{-1})
=\mu(E^{-1})=\mu_R(E)$.
Dla ustalonego $g$ miara $E\mapsto\mu(Eg)$ jest
lewostronnie niezmiennicza, a więc na mocy
twierdzenia~\ref{thm:haar-unique-23} równa $c(g)\mu$.
Złożenie prawych przesunięć daje
$c(gh)=c(g)c(h)$, zaś $c(e)=1$.
Dla niezerowej $f\in C_c(G)$, $f\ge0$,
$c(g)$ jest ilorazem całek
\[
 c(g)=\frac{\int_G f(xg^{-1})\,\mathrm d\mu(x)}
 {\int_G f(x)\,\mathrm d\mu(x)}.
\]
Ciągłość licznika wynika z ciągłości translacji
i lokalnego wspólnego zwartego nośnika funkcji,
zatem $c$ jest ciągła.
\end{proof}

\begin{uwaga}[Dwie konwencje]
W literaturze przez funkcję modularną często oznacza
się $\Delta=c^{-1}$, a nie $c$. Dlatego zawsze
należy sprawdzić wzór definiujący. W tym skrypcie
pozostajemy przy konwencji
$\mu(Eg)=c(g)\mu(E)$.
\end{uwaga}

\begin{wniosek}[Modularność grupy Liego]
\label{cor:haar-adjoint-23}
Niech $G$ będzie skończenie wymiarową grupą Liego,
$\mathfrak g=T_eG$ jej algebrą Liego i
$\operatorname{Ad}_g:\mathfrak g\to\mathfrak g$
różniczką w $e$ sprzężenia $x\mapsto gxg^{-1}$.
Dla konwencji modularnej z~\eqref{eq:haar-modular-23}
zachodzi
\[
 c(g)=|\det(\operatorname{Ad}_{g^{-1}})|
       =|\det(\operatorname{Ad}_{g})|^{-1}.
\]
W szczególności grupa Liego jest unimodularna wtedy
i tylko wtedy, gdy $|\det\operatorname{Ad}_g|=1$
dla każdego $g\in G$.
\end{wniosek}
\begin{proof}
Niech $R_g(x)=xg$, $L_x(y)=xy$.
Przez lewą translację identyfikujemy $T_xG$ i
$T_{xg}G$ z $\mathfrak g$. W tych identyfikacjach
różniczka $R_g$ ma postać
\[
 \mathrm d(L_{(xg)^{-1}})_{xg}\circ
 \mathrm d(R_g)_x\circ\mathrm d(L_x)_e
 =\operatorname{Ad}_{g^{-1}},
\]
co sprawdzamy, podstawiając krzywą $h(t)$ wychodzącą
z $e$: $(xg)^{-1}(xh(t)g)=g^{-1}h(t)g$.
Lewa gęstość objętości na grupie Liego ma w tej
trywializacji stały współczynnik. Prawa translacja
mnoży więc jej lokalną objętość przez
$|\det\operatorname{Ad}_{g^{-1}}|$.
Podstawienie w całce daje
$\mu(Eg)=|\det\operatorname{Ad}_{g^{-1}}|\mu(E)$,
a jednoznaczność czynnika modularnego kończy dowód.
\end{proof}

\begin{przyklad}[Grupy zwarte i dyskretne]
\label{ex:haar-compact-discrete-23}
Na grupie zwartej $G$ normalizujemy miarę Haara
przez $\mu(G)=1$. Dla każdego $g$ równość
$\mu(Gg)=c(g)\mu(G)$ daje $c(g)=1$.
Na grupie dyskretnej miarą Haara jest miara licząca,
więc także $c\equiv1$. Jeżeli grupa jest skończona,
można ją unormować do prawdopodobieństwa przez
podzielenie miary liczącej przez $|G|$.
\end{przyklad}

\begin{przyklad}[$\mathbb R^d$, torus i grupy macierzowe]
\label{ex:haar-standard-23}
Na $(\mathbb R^d,+)$ lewostronną i prawostronną miarą
Haara jest miara Lebesgue'a $\lambda_d$.
Na torusie $\mathbb T^d=\mathbb R^d/\mathbb Z^d$
miarę Haara dostajemy przez przeniesienie
$\lambda_d|_{[0,1)^d}$ na iloraz; jest to miara
probabilistyczna. Na grupie
$\operatorname{GL}_d(\mathbb R)$ lewa i prawa miara Haara
mają gęstość $|\det A|^{-d}$ względem miary
Lebesgue'a na przestrzeni $M_d(\mathbb R)\cong\mathbb R^{d^2}$.
Rzeczywiście, mnożenie przez $B$ z lewej zmienia
jakobian przez $|\det B|^d$, a wyznacznik $A$
przez czynnik $\det B$; oba czynniki się skracają.
Analogicznie dla mnożenia z prawej.
\end{przyklad}

\begin{przyklad}[Nieunimodularna grupa afiniczna]
\label{ex:haar-affine-23}
Rozważmy grupę $G=\mathbb R\rtimes(0,\infty)$
z mnożeniem
$(b,a)(c,t)=(b+ac,at)$. Lewą miarą Haara jest
\[
 \mathrm d\mu(b,a)=\frac{\mathrm d\lambda_2(b,a)}{a^2},
\]
a prawą miarą Haara
$\mathrm d\mu_R(b,a)=a^{-1}\mathrm d\lambda_2(b,a)$.
Przy lewym przesunięciu przez $(c,t)$ jakobian
wynosi $t^2$, a czynnik $a^{-2}$ zmienia się
o $t^{-2}$; otrzymujemy niezmienniczość.
Przy prawym przesunięciu przez $(c,t)$ jakobian
wynosi $t$ i bezpośredni rachunek daje
$\mu(E(c,t))=t^{-1}\mu(E)$.
Jest to prosty przykład $c\not\equiv1$.
\end{przyklad}

\begin{uwaga}[Grupy Liego]
Jeżeli $G$ jest $n$-wymiarową grupą Liego, wybór
niezerowej $n$-formy w $T_eG$ i przeniesienie jej
translacjami lewymi daje lewostronną, nigdzie
nieznikającą formę najwyższego stopnia. Dowodzi to
zarazem, że \emph{każda grupa Liego jest orientowalna},
niezależnie od liczby składowych spójności. Wybierając
orientację zgodną z tą formą i całkując odpowiadającą
jej dodatnią gęstość, otrzymujemy lewą miarę Haara.
W ten sposób konstrukcja analityczna łączy się
z geometrią różniczkową.
\end{uwaga}